%====================================
% KAPITEL 6: FAZIT UND AUSBLICK
%====================================
\section{Fazit und Ausblick}

\subsection{Beantwortung der Forschungsfragen}

Die vorliegende Arbeit verfolgte das Ziel, die Change-Management-Faktoren zu identifizieren, die den Erfolg von ERP-Implementierungen in mittelst\"andischen Unternehmen beeinflussen. Auf Basis einer systematischen Literaturanalyse nach Webster und Watson wurden die zentralen Faktoren extrahiert und mit dem IS-Erfolgsmodell nach DeLone und McLean sowie dem ADKAR-Modell verkn\"upft.

Bezogen auf die \textbf{Hauptforschungsfrage} -- Welche Change-Management-Faktoren beeinflussen den Erfolg von ERP-Implementierungen in mittelst\"andischen Unternehmen? -- lassen sich drei zentrale Faktoren nennen: Kommunikation, F\"uhrungsunterst\"utzung sowie Schulung und Training. Diese Faktoren wirken in unterschiedlicher Intensit\"at auf die Erfolgsdimensionen des DeLone-\&-McLean-Modells. Die Kommunikation adressiert prim\"ar die Nutzerzufriedenheit und Nutzungsintensit\"at, die F\"uhrungsunterst\"utzung den Nettonutzen und die Nutzerzufriedenheit, w\"ahrend Schulung und Training insbesondere die Informationsqualit\"at verbessern.

Die \textbf{Sub-Fragen} wurden schrittweise beantwortet: Die Relevanz des DeLone-\&-McLean-Modells f\"ur den Mittelstand wurde aufgezeigt (Sub-Frage~1), die zentralen Change-Management-Faktoren aus der Literatur ab 2015 identifiziert (Sub-Frage~2), die Rolle von Kommunikation, F\"uhrung und Schulung anhand des ADKAR-Modells bewertet (Sub-Frage~3) und schlie\ss{}lich konkrete Handlungsempfehlungen abgeleitet (Sub-Frage~4). Der entwickelte Bezugsrahmen, der die ADKAR-Phasen mit den D\&-M-Erfolgsdimensionen verkn\"upft, bietet eine strukturierte Grundlage f\"ur die Planung und Bewertung von Change-Management-Ma\ss{}nahmen in ERP-Projekten.

\subsection{Limitationen der Untersuchung}

Die Arbeit unterliegt mehreren methodischen und inhaltlichen Limitationen. Erstens basiert die Untersuchung ausschlie\ss{}lich auf einer Literaturanalyse, nicht auf einer eigenen empirischen Erhebung. Die identifizierten Faktoren und Wirkzusammenh\"ange stellen daher synthetisierte Befunde aus der bestehenden Forschung dar, deren \"Ubertragbarkeit auf den spezifischen Kontext des deutschen Mittelstands einer empirischen Validierung bedarf.

Zweitens wurde der zeitliche Rahmen der Literaturrecherche auf Publikationen ab dem Jahr 2015 beschr\"ankt. W\"ahrend diese Eingrenzung die Aktualit\"at der Befunde sicherstellt, k\"onnten relevante \"altere Beitr\"age, die grundlegende Zusammenh\"ange zwischen Change Management und ERP-Erfolg etabliert haben, nicht vollst\"andig ber\"ucksichtigt werden.

Drittens ist die Heterogenit\"at des Mittelstands als Untersuchungsgegenstand zu nennen. Die Spannweite reicht von Kleinstunternehmen mit wenigen Mitarbeitern bis zu mittleren Unternehmen mit mehreren hundert Besch\"aftigten, sodass generalisierte Aussagen die spezifischen Gegebenheiten des Einzelfalls nicht abbilden k\"onnen. Zudem unterliegen die identifizierten Studien einem m\"oglichen Publication Bias, da positive Ergebnisse h\"aufiger publiziert werden als negative oder nicht signifikante Befunde.

\subsection{Implikationen f\"ur Forschung und Praxis sowie Ausblick}

F\"ur die \textbf{Praxis} bieten die Ergebnisse einen strukturierten Orientierungsrahmen. Unternehmen, die eine ERP-Einf\"uhrung planen, k\"onnen die identifizierten Faktoren als Checkliste nutzen und gezielt Change-Management-Ma\ss{}nahmen entlang der ADKAR-Phasen planen. Besonders die Erkenntnis, dass Kommunikation und F\"uhrungsunterst\"utzung in mittelst\"andischen Strukturen aufgrund der flachen Hierarchien besonders wirkungsvoll eingesetzt werden k\"onnen, ist praxisrelevant.

F\"ur die \textbf{Forschung} ergeben sich mehrere Anschlussfragen. Eine empirische \"Uberpr\"ufung des entwickelten Bezugsrahmens, etwa durch eine quantitative Befragung von ERP-Anwendern im Mittelstand oder durch vergleichende Fallstudien, w\"are ein naheliegender n\"achster Schritt. Dar\"uber hinaus k\"onnte die differenzierte Betrachtung nach Unternehmensgr\"o\ss{}e, Branche oder Implementierungsmethode (traditionell vs. agil) vertiefende Erkenntnisse liefern.

Zusammenfassend l\"asst sich festhalten, dass Change Management ein entscheidender Erfolgsfaktor f\"ur ERP-Implementierungen im Mittelstand ist. Die systematische Ber\"ucksichtigung der identifizierten Faktoren Kommunikation, F\"uhrungsunterst\"utzung sowie Schulung und Training kann dazu beitragen, die hohen Misserfolgsraten von ERP-Projekten zu reduzieren und die Potenziale der Digitalisierung im Mittelstand besser auszusch\"opfen.
